\section{Evaluation}
\label{sec:eval}

\subsection{Method}
All numbers come from the result records of the campaigns (\file{results/}, git-ignored); the plot data in
\file{docs/report/data/} is extracted from them by \file{docs/report/tools/extract\_data.py}. Unless noted:
\begin{itemize}
  \item \textbf{GFLOP/s} means \(2NKM\) divided by the kernel time, i.e.\ cycles counted on the card divided by the
        accel clock measured on the card. End-to-end times include the uploads and the download.
  \item Every row is verified (\cref{sec:host}); all 206 rows of the twelve campaigns passed. The memory-map campaign and the first build's
        bring-up checked every element; the other campaigns used the sampled check (4098 elements per run).
  \item Repetitions of the same run agree to four or five significant digits (the kernel and the DDR controller are
        deterministic for a given size), so the plots show single values. Rebuilding and rerunning the best 200\,MHz
        configuration reproduced \gflops{98.59} exactly.
  \item A build that misses timing on any clock is not measured.
\end{itemize}

\subsection{Overview}
\Cref{fig:progress} shows the best result after each step of the exploration against the peak of that configuration.
Two kinds of steps alternate: \emph{adding compute} (more PEs, a higher clock) raises the peak, and \emph{removing
waiting} (transposed A, larger tiles) brings the measurement closer to it. \Cref{fig:stalls} shows the second kind
through the read-wait counters: the share of kernel cycles in which a read channel was ready but had no data. It is an
upper bound on memory stalls (the kernel can often go on from buffered data), but it tracks the efficiency closely.

\begin{figure}[ht]
  \centering
  % Best GFLOP/s after each step of the exploration, against the peak of that configuration.
\begin{tikzpicture}
\begin{axis}[rpt, ybar=0pt, bar width=14pt, height=6.4cm, ymin=0, ymax=125,
    ylabel={GFLOP/s (kernel)}, xtick=data, x tick label style={rotate=28, anchor=north east, font=\sffamily\scriptsize},
    xticklabels={{PE 8\x8, tile 128}, {PE 8\x8, tile 256}, {PE 16\x8, tile 256}, {PE 32\x8, tile 256},
                 {+ transposed A}, {+ tile 512\x512}, {+ tile 1024\x512}, {+ 225\,MHz}},
    nodes near coords, nodes near coords style={font=\sffamily\scriptsize, /pgf/number format/.cd, fixed, precision=1},
    legend style={at={(0.02,0.98)}, anchor=north west}]
  \addplot[s7, fill=black!8, draw=black!45, bar shift=0pt, nodes near coords=\empty] table[x=step, y=peak] {data/progress.csv};
  \addlegendentry{peak of the configuration}
  \addplot[s1, bar shift=0pt, bar width=9pt] table[x=step, y=gflops] {data/progress.csv};
  \addlegendentry{best measured}
\end{axis}
\end{tikzpicture}

  \caption{Best kernel GFLOP/s after each exploration step (over all sizes measured for that build), with the peak
    \(2P_NP_Mf\) of the configuration.}
  \label{fig:progress}
\end{figure}

\begin{figure}[ht]
  \centering
  % Read-wait shares across the exploration steps, at each build's largest size.
\begin{tikzpicture}
\begin{axis}[rbar, bar width=6pt, height=5.4cm, ymax=100, ylabel={\% of kernel cycles},
    xtick=data, x tick label style={rotate=28, anchor=north east, font=\sffamily\scriptsize},
    xticklabels={{PE 8\x8, tile 128}, {PE 8\x8, tile 256}, {PE 16\x8, tile 256}, {PE 32\x8, tile 256},
                 {+ transposed A}, {+ tile 512\x512}, {+ tile 1024\x512}, {+ 225\,MHz}},
    legend style={at={(0.98,0.98)}, anchor=north east}]
  \addplot[s4] table[x expr=\coordindex, y=await] {data/stalls.csv}; \addlegendentry{A read channel waiting for data}
  \addplot[s2] table[x expr=\coordindex, y=bwait] {data/stalls.csv}; \addlegendentry{B read channel waiting for data}
\end{axis}
\end{tikzpicture}

  \caption{Read waits along the exploration, at each build's largest measured size. The A path dominated until A was
    sent transposed; larger tiles then cut both A and B waits.}
  \label{fig:stalls}
\end{figure}

\subsection{Processing elements}
\Cref{fig:pe_sweep} (tile 256\x256, 200\,MHz): PE 8\x8 reaches \SI{96}{\percent} of its peak and PE 16\x8
\SI{92}{\percent}, but PE 32\x8 only \SI{84.5}{\percent} at \(1024^3\) and \SI{80.4}{\percent} at \(2048^3\). The
\(1000\times1008\times1040\) shape is lower for all three because the kernel computes whole tiles
(\cref{sec:eval:sizes}). For PE 32\x8, the A read channel waited in \SI{93}{\percent} of the kernel cycles: the
on-chip transpose of A reads single 64-byte words, and with four times the PEs of PE 8\x8 the kernel needs A four times
as fast.

\begin{figure}[ht]
  \centering
  % PE sweep: % of peak per size, tile 256, 200 MHz.
\begin{tikzpicture}
\begin{axis}[rbar, bar width=9pt, height=5.6cm, ymin=0, ymax=114, ylabel={\% of peak},
    xtick=data, xticklabels={512\(^3\), 1000\x1008\x1040, 1024\(^3\), 2048\(^3\)},
    nodes near coords, nodes near coords style={font=\sffamily\tiny, rotate=90, anchor=west, /pgf/number format/.cd, fixed, precision=1},
    legend style={at={(0.5,1.02)}, anchor=south, legend columns=3}]
  \addplot[s1] table[x expr=\coordindex, y=pe8pct]  {data/pe_sweep.csv}; \addlegendentry{PE 8\x8 (25.6 GFLOP/s peak)}
  \addplot[s2] table[x expr=\coordindex, y=pe16pct] {data/pe_sweep.csv}; \addlegendentry{PE 16\x8 (51.2)}
  \addplot[s3] table[x expr=\coordindex, y=pe32pct] {data/pe_sweep.csv}; \addlegendentry{PE 32\x8 (102.4)}
\end{axis}
\end{tikzpicture}

  \caption{PE sweep, tile 256\x256, 200\,MHz: percentage of the peak per matrix size.}
  \label{fig:pe_sweep}
\end{figure}

\subsection{Data layout: transposed A}
\label{sec:eval:layout}
With \code{transposed\_a} the host sends \(A^\mathsf{T}\), and the kernel reads A in 16-beat bursts: at \(2048^3\)
the number of A read bursts drops from 2\,097\,152 to 131\,072. \Cref{fig:transposed}: PE 32\x8 goes from 82.3 to
\gflops{87.5} and its A waits from 93 to \SI{16}{\percent}, the same level as B. PE 16\x8 does not change (47.19): its
A waits fell from 53 to \SI{9}{\percent}, but they were mostly hidden by buffering. The transpose costs nothing on the
card and is done once on the host, so transposed A is used for all later PE 32\x8 builds.

\begin{figure}[ht]
  \centering
  % Transposed A: GFLOP/s and the share of cycles the A read channel waits for data, at 2048^3, tile 256.
\begin{tikzpicture}
\begin{groupplot}[group style={group size=2 by 1, horizontal sep=1.6cm}, rbar, ybar=6pt, width=0.5\linewidth, height=5.2cm,
    /pgf/bar width=12pt, xtick=data, xticklabels={PE 16\x8, PE 32\x8}, enlarge x limits=0.5,
    nodes near coords, nodes near coords style={font=\sffamily\scriptsize, /pgf/number format/.cd, fixed, precision=1}]
  \nextgroupplot[ylabel={GFLOP/s at 2048\(^3\)}, ymax=100, legend style={at={(0.02,0.98)}, anchor=north west}]
    \addplot[s7] table[x expr=\coordindex, y=g_plain] {data/transposed.csv}; \addlegendentry{A as is}
    \addplot[s1] table[x expr=\coordindex, y=g_ta]    {data/transposed.csv}; \addlegendentry{A\(^\mathsf{T}\) from the host}
  \nextgroupplot[ylabel={A read waits, \% of kernel cycles}, ymax=105]
    \addplot[s7] table[x expr=\coordindex, y=await_plain] {data/transposed.csv};
    \addplot[s1] table[x expr=\coordindex, y=await_ta]    {data/transposed.csv};
\end{groupplot}
\end{tikzpicture}

  \caption{Transposed A at \(2048^3\), tile 256\x256: GFLOP/s (left) and A read waits (right).}
  \label{fig:transposed}
\end{figure}

\subsection{Memory tiles}
With A and B waiting equally, the remaining loss is the per-tile work: every memory tile reloads its strips of A and
B and drains its C values. Larger tiles mean fewer tiles (\cref{fig:tiling}). \Cref{fig:tiles}: from 256\x256 to
1024\x512, PE 32\x8 goes from 88.3 to \gflops{98.2} at \(4096^3\) (\gflops{98.6} at \(8192^3\), \SI{96.3}{\percent}),
and the waits fall from \SI{16.6}{\percent} to \SI{3.8}{\percent} (A) and \SI{2.8}{\percent} (B). The N tile matters
more than the M tile: 512\x256 reaches 94.4 but 256\x512 only 89.4, because a larger N tile means fewer passes over B.

The cost is BRAM only: LUTs and DSPs stay the same, while the kernel's BRAM36 count goes from 285 to 509 and the
design from 414 to 654.5 BRAM tiles (\SI{69}{\percent}). A 1024\x1024 tile needs 2070 BRAM18 by the HLS estimate,
more than the 1910 on the device, so 1024\x512 is the largest tile for PE 32\x8.

\begin{figure}[ht]
  \centering
  % Memory tile sweep for PE 32x8 with transposed A, 200 MHz.
\begin{tikzpicture}
\begin{groupplot}[group style={group size=2 by 1, horizontal sep=1.7cm}, rpt, width=0.47\linewidth, height=5.6cm,
    xtick=data, xticklabels={256\x256, 256\x512, 512\x256, 512\x512, 1024\x512},
    x tick label style={rotate=25, anchor=north east}, xlabel={memory tile \(T_N\times T_M\)}]
  \nextgroupplot[ybar, bar width=5pt, ymin=80, ymax=105, ylabel={GFLOP/s}, enlarge x limits=0.14,
      legend style={at={(0.02,0.98)}, anchor=north west}]
    \addplot[s6] table[x expr=\coordindex, y=g2048] {data/tiles.csv}; \addlegendentry{2048\(^3\)}
    \addplot[s1] table[x expr=\coordindex, y=g4096] {data/tiles.csv}; \addlegendentry{4096\(^3\)}
    \addplot[s5] table[x expr=\coordindex, y=g8192] {data/tiles.csv}; \addlegendentry{8192\(^3\)}
    \draw[densely dashed, cRed] (axis cs:-0.5,102.4) -- node[pos=0.5, above, font=\sffamily\tiny, text=cRed] {peak 102.4} (axis cs:4.5,102.4);
  \nextgroupplot[ymin=0, ymax=20, ylabel={read waits at 4096\(^3\), \% of cycles}, enlarge x limits=0.1,
      axis y line*=left, legend style={at={(0.98,0.98)}, anchor=north east}]
    \addplot[l4, mark=*] table[x expr=\coordindex, y=await4096] {data/tiles.csv}; \addlegendentry{A waits}
    \addplot[l2, mark=square*] table[x expr=\coordindex, y=bwait4096] {data/tiles.csv}; \addlegendentry{B waits}
\end{groupplot}
\end{tikzpicture}

  \caption{Memory tile sweep for PE 32\x8 with transposed A, 200\,MHz. Left: GFLOP/s at three sizes (8192\(^3\) was
    measured for the two largest tiles). Right: read waits at \(4096^3\).}
  \label{fig:tiles}
\end{figure}

\subsection{Matrix sizes}
\label{sec:eval:sizes}
The size sweep ran 45 shapes on the best 200\,MHz build (\cref{fig:size}). Two effects explain all of it:
\begin{itemize}
  \item \textbf{Whole tiles.} With \(T_N = 1024\), a matrix with \(N = 256\) runs at a quarter of the speed of one
        with \(N \ge 1024\); with \(T_M = 512\), \(M = 256\) halves it. \(N = M = K = 256\) reaches only 11
        GFLOP/s. For small matrices a smaller-tile build is faster: the 256\x256 build reaches 84 GFLOP/s at
        \(512^3\), while the 1024\x512 build reaches 48 for \(N = 512\) even with \(K = M = 1024\).
  \item \textbf{Fill and drain per tile}, amortised over \(K\): with \(K=256\) the kernel reaches 87.5, with
        \(K = 4096\) 98.2 GFLOP/s.
\end{itemize}
End to end (dashed), the transfers cost a lot below \(\sim\)\(1024^3\), where the kernel takes only milliseconds; at
\(4096^3\) and above the kernel dominates (\cref{sec:eval:e2e}).

\begin{figure}[ht]
  \centering
  \begin{subfigure}[b]{\linewidth}\centering% Size sweep on the best 200 MHz build (tile 1024x512): kernel GFLOP/s against N for each (K, M).
\begin{tikzpicture}
\begin{axis}[rpt, height=6.3cm, xmode=log, log basis x=2, xtick={256,512,1024,2048,4096},
    xticklabels={256,512,1024,2048,4096}, xlabel={\(N\)}, ylabel={GFLOP/s (kernel)}, ymin=0, ymax=105,
    enlarge x limits=0.04, legend columns=3, legend style={at={(0.5,1.02)}, anchor=south}]
  \addplot[l1, mark=*]        table[x=n, y=k4096m4096] {data/size_kernel.csv}; \addlegendentry{\(K\)=4096, \(M\)=4096}
  \addplot[l3, mark=square*]  table[x=n, y=k1024m1024] {data/size_kernel.csv}; \addlegendentry{\(K\)=1024, \(M\)=1024}
  \addplot[l2, mark=triangle*]table[x=n, y=k256m1024]  {data/size_kernel.csv}; \addlegendentry{\(K\)=256, \(M\)=1024}
  \addplot[l5, mark=diamond*] table[x=n, y=k4096m256]  {data/size_kernel.csv}; \addlegendentry{\(K\)=4096, \(M\)=256}
  \addplot[l4, mark=pentagon*]table[x=n, y=k256m256]   {data/size_kernel.csv}; \addlegendentry{\(K\)=256, \(M\)=256}
  \addplot[l1, densely dashed, mark=o] table[x=n, y=k4096m4096] {data/size_e2e.csv}; \addlegendentry{end to end, 4096/4096}
  \draw[black!40, densely dotted] (axis cs:1024,0) -- (axis cs:1024,105) node[pos=0.04, right, font=\sffamily\tiny, text=black!55] {\(N = T_N\)};
  \draw[densely dashed, cRed] (axis cs:256,102.4) -- (axis cs:4096,102.4) node[pos=0.03, above, font=\sffamily\tiny, text=cRed, anchor=south west] {peak 102.4};
\end{axis}
\end{tikzpicture}
\end{subfigure}\\[3mm]
  \begin{subfigure}[b]{\linewidth}\centering% Kernel GFLOP/s over N and M at K = 4096 (tile 1024x512).
\begin{tikzpicture}
\begin{axis}[width=0.62\linewidth, height=4.4cm, view={0}{90}, colormap/viridis, colorbar,
    colorbar style={width=2.5mm, ylabel={GFLOP/s}, font=\sffamily\scriptsize},
    xtick={0,...,4}, xticklabels={256,512,1024,2048,4096}, ytick={0,1,2}, yticklabels={256,1024,4096},
    xlabel={\(N\)}, ylabel={\(M\)}, xmin=-0.5, xmax=4.5, ymin=-0.5, ymax=2.5, axis on top,
    nodes near coords={\pgfmathprintnumber[fixed, precision=1]{\pgfplotspointmeta}},
    nodes near coords style={font=\sffamily\scriptsize, text=black, fill=white, fill opacity=0.65, text opacity=1, inner sep=1pt, anchor=center}]
  \addplot[matrix plot*, mesh/rows=3, mesh/ordering=y varies, point meta=explicit]
    table[x=xi, y=yi, meta=gflops] {data/size_heat.csv};
\end{axis}
\end{tikzpicture}
\end{subfigure}
  \caption{Size sweep on PE 32\x8, transposed A, tile 1024\x512, 200\,MHz. Top: kernel GFLOP/s against \(N\) for
    several \((K, M)\), and end to end for \(K = M = 4096\). Bottom: GFLOP/s over \(N\) and \(M\) at \(K = 4096\);
    the steps are the whole-tile effect.}
  \label{fig:size}
\end{figure}

\subsection{Roofline}
\label{sec:eval:roofline}
\Cref{fig:roofline} places every measured build in a roofline model~\cite{roofline}. The intensity is that of one
memory tile (A and B traffic; C is written once), the memory roof one DDR3 channel at the measured \SI{7.2}{GB/s}.
Even the smallest tile (\(128^2\), 32 flop/byte) needs only \SI{0.8}{GB/s} at PE 8\x8's peak, and the best build
needs about \SI{0.65}{GB/s} at \gflops{110.9}. Every point is far to the right of the ridge, so bandwidth is never
the limit; what kept the kernel from its peak was latency and the per-tile overheads, which is why
bursts and larger tiles helped and why the memory map hardly matters (\cref{sec:eval:maps}).

\begin{figure}[ht]
  \centering
  % Roofline: tile intensity (A + B traffic) against measured GFLOP/s, with one DDR3 channel as the memory roof.
\begin{tikzpicture}
\begin{axis}[rpt, height=6.4cm, xmode=log, ymode=log, log basis x=10, log basis y=10, xmin=4, xmax=400, ymin=8, ymax=200,
    xlabel={arithmetic intensity of a memory tile (flop per byte of A + B)}, ylabel={GFLOP/s},
    legend style={at={(0.98,0.02)}, anchor=south east}, log ticks with fixed point]
  \addplot[l7, very thick, domain=4:400, samples=200] {min(7.2*x, 1e9)}; \addlegendentry{one DDR3 channel, 7.2\,GB/s}
  \foreach \p/\c in {25.6/cBlue, 51.2/cOrange, 102.4/cGreen, 115.2/cPurple} {
    \edef\tmp{\noexpand\draw[\c, densely dashed] (axis cs:4,\p) -- (axis cs:400,\p) node[pos=0.0, above right, font=\noexpand\sffamily\noexpand\tiny, text=\c] {\p};}
    \tmp
  }
  \addplot[only marks, mark=*, mark size=2.2pt, cBlue] table[x expr={\thisrow{tn}*\thisrow{tm}/(2*(\thisrow{tn}+\thisrow{tm}))},
      y=gflops, restrict expr to domain={\thisrow{pn}}{8:8}] {data/builds.csv}; \addlegendentry{PE 8\x8}
  \addplot[only marks, mark=square*, mark size=2.2pt, cOrange] table[x expr={\thisrow{tn}*\thisrow{tm}/(2*(\thisrow{tn}+\thisrow{tm}))},
      y=gflops, restrict expr to domain={\thisrow{pn}}{16:16}] {data/builds.csv}; \addlegendentry{PE 16\x8}
  \addplot[only marks, mark=triangle*, mark size=2.8pt, cGreen!70!black] table[x expr={\thisrow{tn}*\thisrow{tm}/(2*(\thisrow{tn}+\thisrow{tm}))},
      y=gflops, restrict expr to domain={\thisrow{pn}}{32:32}] {data/builds.csv}; \addlegendentry{PE 32\x8}
\end{axis}
\end{tikzpicture}

  \caption{Roofline: tile intensity against measured GFLOP/s for every build, with the compute peaks of 25.6, 51.2,
    102.4 and 115.2\,GFLOP/s and one DDR3 channel as the memory roof.}
  \label{fig:roofline}
\end{figure}

\subsection{Clock frequency and Vivado strategies}
\label{sec:eval:timing}
\paragraph{Clock.} The HLS clock estimate (4.24\,ns for PE 8\x8, 4.21\,ns for PE 32\x8) does not change between
200 and 250\,MHz targets, but HLS pipelines deeper for a tighter target: +7.2\,k flip-flops for PE 8\x8 at 250\,MHz,
+18.6\,k for PE 32\x8 at 225\,MHz. Whether that is enough only a routed build shows (\cref{fig:clock}):
\begin{itemize}
  \item \textbf{PE 8\x8 closes at 250\,MHz} (accel slack +0.049\,ns) and reaches \gflops{29.0}. Throughput scales
        almost linearly with the clock; kernel cycles grow by \SI{1.5}{\percent} from 150 to 250\,MHz because the DDR
        controllers stay at 125\,MHz.
  \item \textbf{PE 32\x8 misses 225\,MHz by 41\,ps} with the default strategy (34 endpoints). The failing paths are
        routing: the 256-bit C-pipe FIFOs between neighbouring PEs, whose write enable fans out to 256 shift-register
        bits over 4.3--4.5\,ns of wire behind a single LUT, and the dataflow control's ready/reset fan-out.
  \item \textbf{At 250\,MHz PE 32\x8 misses by 0.78\,ns} (8455 endpoints) on the dataflow control paths (7 logic
        levels). Going there would need kernel changes, e.g.\ registering the control fan-out.
\end{itemize}

\paragraph{Presets.} \Cref{fig:presets}: the presets do not change size or cycles, only closure. At 225\,MHz the
\code{explore} preset closes PE 32\x8 with +0.007\,ns and \code{aggressive} with +0.049\,ns; its physical
optimisation replicates exactly the high-fanout drivers that failed. The result is \textbf{\gflops{110.9} at
\(8192^3\)}, \SI{96.2}{\percent} of the peak, \SI{12.4}{\percent} faster than at 200\,MHz, with the same kernel cycle
count (2.2316 vs 2.2304\,G), so the kernel is still compute-bound at 225\,MHz. Both preset builds give identical
results. \code{aggressive} costs about \SI{9}{\percent} more build time and gave the most margin on every clock in
both experiments, but one build per preset is a hint, not a ranking.

\begin{figure}[ht]
  \centering
  % Accel clock: GFLOP/s (PE 8x8) and routed slack of the accel clock and of the PCIe core clock.
\begin{tikzpicture}
\begin{groupplot}[group style={group size=2 by 1, horizontal sep=1.8cm}, rpt, width=0.5\linewidth, height=5.4cm,
    xlabel={accel clock (MHz)}, xtick={150,200,225,250}]
  \nextgroupplot[ylabel={GFLOP/s at 1024\(^3\) (PE 8\x8)}, ymin=0, ymax=35, xmin=140, xmax=260,
      legend style={at={(0.02,0.98)}, anchor=north west}]
    \addplot[l7, densely dashed, domain=150:250, mark=none] {23.355*x/200}; \addlegendentry{linear in the clock}
    \addplot[l1, mark=*, only marks, mark size=2.5pt, nodes near coords, nodes near coords style={font=\sffamily\scriptsize, anchor=south east}]
      table[x=mhz, y=gflops, restrict expr to domain={\coordindex}{0:2}] {data/clock.csv};
    \addlegendentry{measured}
  \nextgroupplot[ylabel={WNS (ns)}, ymin=-0.85, ymax=0.15, xmin=140, xmax=260,
      legend style={at={(0.02,0.02)}, anchor=south west}]
    \draw[black!60] (axis cs:140,0) -- (axis cs:260,0);
    \addplot[l2, mark=*]       table[x=mhz, y=wns_accel, restrict expr to domain={\coordindex}{3:5}] {data/clock.csv};
    \addlegendentry{PE 32\x8: accel clock}
    \addplot[l1, mark=square*] table[x=mhz, y=wns_accel, restrict expr to domain={\coordindex}{0:2}] {data/clock.csv};
    \addlegendentry{PE 8\x8: accel clock}
    \addplot[only marks, mark=triangle*, mark size=2.6pt, cGreen!70!black] table[x=mhz, y=wns_pcie] {data/clock.csv};
    \addlegendentry{PCIe core 500\,MHz (either)}
\end{groupplot}
\end{tikzpicture}

  \caption{Left: PE 8\x8 GFLOP/s against the accel clock (the 200\,MHz value is from the \code{explore} build; the
    default-preset build missed timing on the PCIe path). Right: routed slack of the accel clock for both
    configurations and of the PCIe core clock, default preset.}
  \label{fig:clock}
\end{figure}

\begin{figure}[ht]
  \centering
  % Vivado presets: routed WNS per clock for the two configurations where it mattered.
\begin{tikzpicture}
\begin{axis}[rbar, ybar=1pt, bar width=6pt, height=5.6cm, ymin=-0.07, ymax=0.25, ylabel={WNS (ns)},
    xtick=data, xticklabels={default, explore, aggressive, default, explore, aggressive},
    extra x ticks={1,4}, extra x tick labels={PE 8\x8 @ 200\,MHz, PE 32\x8 @ 225\,MHz},
    extra x tick style={tick label style={yshift=-4.2mm, font=\sffamily\small}, major tick length=0pt, grid=none},
    legend style={at={(0.5,1.02)}, anchor=south, legend columns=3}, enlarge x limits=0.1]
  \addplot[s2] table[x expr=\coordindex, y=wns_accel] {data/presets.csv}; \addlegendentry{accel}
  \addplot[s3] table[x expr=\coordindex, y=wns_pcie]  {data/presets.csv}; \addlegendentry{PCIe core 500\,MHz}
  \addplot[s1] table[x expr=\coordindex, y=wns_sys]   {data/presets.csv}; \addlegendentry{sys 125\,MHz}
  \draw[black!70] (axis cs:-0.5,0) -- (axis cs:5.5,0);
  \draw[black!30] (axis cs:2.5,-0.07) -- (axis cs:2.5,0.25);
\end{axis}
\end{tikzpicture}

  \caption{Routed slack per clock for the three Vivado presets, in the two experiments.}
  \label{fig:presets}
\end{figure}

\subsection{Timing margin and the PCIe core path}
\Cref{fig:slack} shows the slack of the accel clock and of the PCIe core clock in all 20 builds. The accel clock
failed only where the clock was pushed. The PCIe core clock (500\,MHz in the hard block for Gen2 x8) has between
\(-0.009\) and \(+0.076\)\,ns in every build, whatever the kernel: its critical path goes from the \code{PCIE\_2\_1}
block to its own receive-buffer BRAM, which the Xilinx IP places next to it, with no logic in between (1.375\,ns of
the 2\,ns period). It failed alone once (PE 8\x8, 200\,MHz, default preset); that build was not used. Configuring the
PHY as Gen2 x4, the width the development PC's slot provides, would halve that clock; it is listed as an open decision.

\begin{figure}[ht]
  \centering
  % Routed WNS of the PCIe core clock (500 MHz) and of the accel clock in every build.
\begin{tikzpicture}
\begin{axis}[rpt, ybar=0pt, bar width=3.5pt, height=6.0cm, ylabel={WNS (ns)}, ymin=-0.1, ymax=0.25,
    xtick=data, xticklabels from table={data/builds.csv}{label},
    x tick label style={rotate=60, anchor=north east, font=\sffamily\tiny}, enlarge x limits=0.03,
    legend style={at={(0.98,0.98)}, anchor=north east}, restrict y to domain=-0.1:0.3]
  \addplot[s3] table[x expr=\coordindex, y=wns_pcie]  {data/builds.csv}; \addlegendentry{PCIe core 500\,MHz}
  \addplot[s2] table[x expr=\coordindex, y=wns_accel] {data/builds.csv}; \addlegendentry{accel (bars below \(-0.1\) clipped)}
  \draw[black!70] (axis cs:-0.5,0) -- (axis cs:19.5,0);
\end{axis}
\end{tikzpicture}

  \caption{Routed slack (WNS) of the PCIe core clock and of the accel clock in every build. The PCIe path is
    marginal in all of them; the accel clock fails only at 225\,MHz (default preset) and 250\,MHz for PE 32\x8.}
  \label{fig:slack}
\end{figure}

\subsection{Resources}
\Cref{fig:resources} splits the utilisation of the best build into the kernel and the rest of the SoC. DSPs and
BRAM are the scarce resources: each multiply-add uses five DSP48E1 slices (three for the multiplier, two for the adder),
so the 256 MACs of PE 32\x8 need 1280 of the 1320 DSPs in the kernel; and the C tile fills most of the kernel's 543.5
BRAM tiles. A PE 64\x8 array would need about 2600 DSPs with this mapping, more than the 1920 on the device.

\begin{figure}[ht]
  \centering
  % Device utilisation of the best build (PE 32x8, tile 1024x512, 225 MHz): kernel vs the rest of the SoC.
\begin{tikzpicture}
\begin{axis}[rpt, xbar stacked, enlarge x limits=false, width=0.84\linewidth, bar width=9pt, height=4.4cm, xmin=0, xmax=100, xlabel={\% of the xc7k480t},
    ytick={0,1,2,3}, yticklabels={LUT (298.6\,k), FF (597.2\,k), BRAM tiles (955), DSP48 (1920)}, y dir=reverse,
    enlarge y limits=0.2, legend style={at={(0.98,0.02)}, anchor=south east}, grid=none, xmajorgrids]
  % kernel: LUT 122.7k, FF 206.8k, BRAM 509 + 69/2 = 543.5 tiles, DSP 1320; whole design: 152.0k, 235.8k, 654.5, 1320
  \addplot[s2] coordinates {(41.1,0) (34.6,1) (56.9,2) (68.75,3)}; \addlegendentry{GEMM kernel}
  \addplot[s1] coordinates {(9.8,0) (4.9,1) (11.6,2) (0,3)};      \addlegendentry{rest of the SoC}
\end{axis}
\end{tikzpicture}

  \caption{Utilisation of the best build (PE 32\x8, tile 1024\x512, 225\,MHz, \code{aggressive}).}
  \label{fig:resources}
\end{figure}

\subsection{Memory maps}
\label{sec:eval:maps}
Two other maps were built and verified (PE 8\x8, every element checked): all three matrices on CH0, and A on CH1
with B and C on CH0 (\cref{fig:memmap_plot}). Sharing a channel between A and B costs \SI{0.3}{\percent}: B's read
latency rises from 39 to 46 cycles and its waits from 15 to \SI{18}{\percent}. At PE 8\x8 the kernel needs well under
\SI{1}{GB/s}, so this confirms the roofline picture; the map should matter more for PE 32\x8, which was not measured.

\begin{figure}[ht]
  \centering
  % Memory maps for PE 8x8, tile 128, 200 MHz: GFLOP/s and B read latency at 2048^3.
\begin{tikzpicture}
\begin{groupplot}[group style={group size=2 by 1, horizontal sep=1.6cm}, rbar, ybar=4pt, width=0.5\linewidth, height=4.6cm,
    /pgf/bar width=13pt, xtick=data, xticklabels={{A0 B1 C1}, {A0 B0 C0}, {A1 B0 C0}}, enlarge x limits=0.3,
    nodes near coords, nodes near coords style={font=\sffamily\scriptsize, /pgf/number format/.cd, fixed, precision=2}]
  \nextgroupplot[ylabel={GFLOP/s at 2048\(^3\)}, ymin=23.2, ymax=23.5]
    \addplot[s1] table[x expr=\coordindex, y=g2048] {data/memmap.csv};
  \nextgroupplot[ylabel={B read latency (cycles)}, ymin=0, ymax=55]
    \addplot[s2] table[x expr=\coordindex, y=blat2048] {data/memmap.csv};
\end{groupplot}
\end{tikzpicture}

  \caption{Memory maps (A, B, C on channel 0 or 1), PE 8\x8, tile 128, 200\,MHz, at \(2048^3\).}
  \label{fig:memmap_plot}
\end{figure}

\subsection{Transfers and end-to-end time}
\label{sec:eval:e2e}
At \(8192^3\) the host uploads 512\,MB (A and B) and downloads 256\,MB (C). With \SI{0.82}{GB/s} up and
\SI{1.49}{GB/s} down (\cref{sec:pcie:rates}) that takes 0.83\,s next to 9.92\,s of kernel time: the kernel is
\SI{92}{\percent} of the 10.77\,s end to end (\cref{fig:e2e}). At \(2048^3\) it is only \SI{74}{\percent}, and the
upload alone is \SI{19}{\percent}. The H2C stop delay (\cref{sec:pcie:rates}) is therefore the next thing to fix for
mid-size matrices.

\begin{figure}[ht]
  \centering
  % End-to-end time of one GEMM on the best build, split into kernel and transfers.
\begin{tikzpicture}
\begin{axis}[rpt, xbar stacked, enlarge x limits=false, width=0.84\linewidth, bar width=10pt, height=3.8cm, xmin=0, xmax=100, xlabel={\% of end-to-end time},
    ytick={0,1,2}, yticklabels={2048\(^3\) (0.21\,s), 4096\(^3\) (1.46\,s), 8192\(^3\) (10.77\,s)}, y dir=reverse,
    enlarge y limits=0.3, legend style={at={(0.5,1.03)}, anchor=south, legend columns=4}, grid=none, xmajorgrids]
  \addplot[s2] table[y expr=\coordindex, x=kernel_pct] {data/e2e.csv}; \addlegendentry{kernel}
  \addplot[s1] table[y expr=\coordindex, x=h2c_pct]    {data/e2e.csv}; \addlegendentry{upload A, B}
  \addplot[s3] table[y expr=\coordindex, x=c2h_pct]    {data/e2e.csv}; \addlegendentry{download C}
  \addplot[s7] table[y expr=\coordindex, x=other_pct]  {data/e2e.csv}; \addlegendentry{other}
\end{axis}
\end{tikzpicture}

  \caption{Where the end-to-end time goes on the best build.}
  \label{fig:e2e}
\end{figure}
